\ysrc{203}
定理的第二部分相当形式化；第一部分由 [FF18, Thm. 9.3.1] 推出。

\begin{rem}{5.1.2}\label{rem:a4-5-1-2}
在 Scholze 的菱形（diamond）世界里，有 $X_{F_\infty}/\Gamma=\overline{X}_F/\mathrm{Gal}_F$，这给第一个等价一个更几何的解释。
\end{rem}

此外，这些等价之间有显然的相容性。借助 Beauville--Laszlo 粘合引理 [BL95]，还可以给出这些范畴的另一种描述，由此得到 Berger 的 $B$-对（$B$-pairs）范畴 [Ber08a]。我们置
\[
B_e=\Gamma(\overline{X}_F\setminus\{\infty\},\mathcal{O}_{\overline{X}_F}),\qquad
B^+_{\mathrm{dR}}=\widehat{\mathcal{O}}_{\overline{X}_F,\infty},\qquad
B_{\mathrm{dR}}=B^+_{\mathrm{dR}}[1/t].
\]
这些环中的每一个都带有连续的 $\mathrm{Gal}_F$-作用。利用 Beauville--Laszlo 粘合引理，可把 $\mathrm{VB}^{\mathrm{Gal}_F}_{\overline{X}_F}$ 等同于三元组 $(M,\Lambda,\xi)$ 的范畴：其中 $M$ 是有限投射 $B_e$-模，$\Lambda$ 是有限投射 $B^+_{\mathrm{dR}}$-模，两者都带有连续半线性的 $\mathrm{Gal}_F$-作用，而 $\xi$ 是 $\mathrm{Gal}_F$-等变的同构
\[
\xi:M\otimes_{B_e}B_{\mathrm{dR}}\xrightarrow{\ \sim\ }\Lambda\otimes_{B^+_{\mathrm{dR}}}B_{\mathrm{dR}}.
\]
对 $\mathrm{VB}^{\mathrm{Gal}_F}_{\overline{X}_F}\cong\mathrm{VB}^{\Gamma}_{X_{F_\infty}}$ 还有第三种描述，用非完美周期环（imperfect period ring）上的 $(\varphi,\Gamma)$-模给出。记 $\mathcal{R}_F$ 为 Robba 环
\[
\mathcal{R}_F=\varinjlim_{r\to1}\varprojlim_{s\to1}B^{[r,s]}_F,
\]
其中 $B^{[r,s]}_F$ 是闭圆环 $\mathbf{B}^{[r,s]}_{F_{\infty,0}}$ 上刚性解析函数构成的环，该圆环位于 $F_{\infty,0}$ 上的闭单位圆盘内，内半径为 $r$、外半径为 $s$；这里 $F_{\infty,0}$ 是 $F_\infty$ 的极大非分歧子扩张（它一般可以严格大于 $F_0$）。

当 $s>r\gg0$ 时，有典范的方式在 $B^{[r,s]}_F$ 上定义连续的 $\Gamma$-作用，并定义与 $\Gamma$-作用交换的环同态
\[
\varphi:B^{[r,s]}_F\to B^{[r^{1/p},s^{1/p}]}_F.
\]
这些数据诱导出环同态 $\varphi:\mathcal{R}_F\to\mathcal{R}_F$，以及 $\mathcal{R}_F$ 上与 $\varphi$ 交换的连续 $\Gamma$-作用（$\varphi$-作用与 $\Gamma$-作用的定义细节，以及下文出现的、系数取仿射代数的 Robba 环 $\mathcal{R}_F$ 的各种变体，可参看 [KPX14, 2.2]）。

\begin{defn}{5.1.3}\label{defn:a4-5-1-3}
范畴 $\Mod^{\varphi,\Gamma}_{\mathcal{R}_F}$ 由有限投射 $\mathcal{R}_F$-模、其上的连续 $\Gamma$-作用，以及与 $\Gamma$-作用交换的同构 $\varphi^*D\to D$ 组成。
\end{defn}

\begin{rem}{5.1.4}\label{rem:a4-5-1-4}
我们回忆 $\mathcal{R}_F$ 上的 $(\varphi,\Gamma)$-模与第 4 节所用 $(\varphi,\Gamma)$-模的关系。第 4 节的 $(\varphi,\Gamma)$-模定义在一个 $p$-进完备环 $\mathbf{A}_F$ 上。该环含有一个由超收敛元素构成的子环 $\mathbf{A}^\dagger_F\subset\mathbf{A}_F$，它在 $\varphi$ 与 $\Gamma$ 的作用下稳定。此外，$\mathbf{A}^\dagger_F$ 等变地（对 $\varphi$ 与 $\Gamma$）嵌入 $\Q_p$-代数 $\mathcal{R}_F$。由 Cherbonnier--Colmez 的一个定理 [CC98]，$\mathbf{B}_F=\mathbf{A}_F[1/p]$ 上每个 étale $(\varphi,\Gamma)$-模都是超收敛的；按定义，$\mathbf{B}_F$ 上的 étale $(\varphi,\Gamma)$-模就是那些由 $\mathbf{A}_F$ 上的 $(\varphi,\Gamma)$-模经基变换得到的模。更确切地说，从 $\mathbf{B}^\dagger_F=\mathbf{A}^\dagger_F[1/p]$ 到 $\mathbf{B}_F$ 的标量扩张，在相应的 étale $(\varphi,\Gamma)$-模范畴之间诱导出等价。\ysrc{204}此外，从 $\mathbf{B}^\dagger_F$ 到 $\mathcal{R}_F$ 的标量扩张，把 $\mathbf{B}^\dagger_F$ 上的 étale $(\varphi,\Gamma)$-模范畴全忠实地嵌入 $\mathcal{R}_F$ 上的 $(\varphi,\Gamma)$-模范畴。按定义，该函子本质像中的 $\mathcal{R}_F$ 上的 $(\varphi,\Gamma)$-模称为 étale 的。
\end{rem}

由 Kedlaya 的斜率滤链定理 [Ked08]，étale $(\varphi,\Gamma)$-模恰为斜率为零的半稳定对象，含义如下：设 $D$ 是秩为 $r$ 的 $(\varphi,\Gamma)$-模\tn{底本此句作「If $D$ is a $(\varphi,\Gamma)$-module $D$ of rank $r$」，其中 $D$ 重复出现，译文按语义只保留一处。}，则由秩 $1$ 的 $(\varphi,\Gamma)$-模的分类（见下文 §7.1），可写 $\bigwedge^rD=\mathcal{R}_F(\delta)$，其中 $\delta$ 是 $F^\times$ 的连续特征。于是定义 $D$ 的斜率为
\[
\mathrm{slope}(D)=\frac{\mathrm{val}_p(\delta(\varpi_F))}{r},
\]
其中 $\varpi_F$ 是 $F$ 中单值化元的一个选取。进而称 $D$ 半稳定，若对一切 $(\varphi,\Gamma)$-稳定子对象 $D'\subseteq D$ 都有 $\mathrm{slope}(D')\ge\mathrm{slope}(D)$。

\begin{thm}{5.1.5}\label{thm:a4-5-1-5}
存在范畴的等价
\[
\Mod^{\varphi,\Gamma}_{\mathcal{R}_F}\longrightarrow\mathrm{VB}^{\varphi,\Gamma}_{Y_{F_\infty}}.
\]
\end{thm}

\begin{proof}[证明梗概]
该定理背后的粗略想法如下：$\mathcal{R}_F$ 上任一 $(\varphi,\Gamma)$-模都在某个闭圆环 $\mathbf{B}^{[r,s]}_{F_{\infty,0}}$ 上有模型。利用坐标函数 $[\varpi^\flat]$ 的选取，可把 $Y_{F_\infty}$ 中相应的闭圆环映到 $\mathbf{B}^{[r,s]}_{F_{\infty,0}}$，且该映射对 $\Gamma$-作用是等变的（并以显然的方式与 $\varphi$ 相容）。把各个 $\mathbf{B}^{[r,s]}_{F_{\infty,0}}$ 上带 $\Gamma$-作用的 $\varphi$-丛拉回到 $Y_{F_\infty}$ 中的闭圆环，再将这些拉回扩张，便得到 $Y_{F_\infty}$ 上 $\Gamma$-等变的 $\varphi$-丛（因为 $\varphi$ 是 $Y_{F_\infty}$ 的自同构，它会移动闭圆环的半径）。定理的困难部分是本质满射性（有时称为去完美化，deperfection），其证明用到分圆塔上由 Tate 正规化迹给出的额外结构。定理的证明可由 [FF18, 10.1.2] 与 [Ber08a, Theorem 2.2.7] 合并得到；用 Tate 正规化迹作下降论证的变体，可参看 [BC08, 4.2]。
\end{proof}

我们关心的是这些（彼此等价的）对象以刚性解析空间为系数的族。更确切地说，我们想定义一个叠 $\mathfrak{X}_d$，其 $\Sp(\Q_p)$-值点由 $\mathrm{VB}^{\mathrm{Gal}_F}_{\overline{X}_F}$ 中秩为 $d$ 的对象（或上述任一等价范畴中秩为 $d$ 的对象）构成的群胚给出。下面简述上述范畴在仿射刚性空间 $\Sp(A)$ 上的变体，这已足以定义我们的叠。

设 $A$ 是 $\Q_p$ 上拓扑有限型的仿射代数。由 Kedlaya 的一个结果 [Ked16]，Fargues--Fontaine 曲线 $X_{F_\infty}$、$\overline{X}_F$ 都是强诺特的，它们的覆盖 $Y_{F_\infty}$、$\overline{Y}_F$ 也（局部地）如此。因此纤维积
\[
X_{F_\infty,A}:=X_{F_\infty}\times_{\Sp(\Q_p)}\Sp(A)
\]
等等在 adic 空间范畴中有定义，并带有 $\Gamma$（相应地 $\mathrm{Gal}_F$）的典范连续作用。

与上面考虑的绝对情形类似，可以定义
\begin{itemize}[leftmargin=1.8em,itemsep=0.15em]
\item $\Gamma$-等变向量丛在 $X_{F_\infty,A}$ 上的范畴 $\mathrm{VB}^{\Gamma}_{X_{F_\infty,A}}$，
\item $\mathrm{Gal}_F$-等变向量丛在 $\overline{X}_{F,A}$ 上的范畴 $\mathrm{VB}^{\mathrm{Gal}_F}_{\overline{X}_{F,A}}$，
\item $(\varphi,\Gamma)$-等变向量丛在 $Y_{F_\infty,A}$ 上的范畴 $\mathrm{VB}^{\varphi,\Gamma}_{Y_{F_\infty,A}}$，
\item $(\varphi,\mathrm{Gal}_F)$-等变向量丛在 $\overline{Y}_{F,A}$ 上的范畴 $\mathrm{VB}^{\varphi,\mathrm{Gal}_F}_{\overline{Y}_{F,A}}$。
\end{itemize}

\ysrc{205}
再次指出，向量丛上的 $\Gamma$-（相应地 $\mathrm{Gal}_F$-）作用应当按显然的方式理解成连续的。此外，$(\varphi,\Gamma)$-模的范畴还有以非完美 Robba 环 $\mathcal{R}_F$ 为系数的变体。定义
\[
\mathcal{R}_{F,A}=\varinjlim_{r\to1}\varprojlim_{s\to1}B^{[r,s]}_F\widehat{\otimes}_{\Q_p}A,
\]
它同样带有连续 $\Gamma$-作用以及与 $\Gamma$-作用交换的自同态 $\varphi$。于是范畴 $\Mod^{\varphi,\Gamma}_{\mathcal{R}_{F,A}}$ 定义为：对象是有限投射 $\mathcal{R}_{F,A}$-模 $D$，配有同构 $\varphi_D:\varphi^*D\to D$ 以及与 $\varphi_D$ 交换的连续 $\Gamma$-作用。

定理 5.1.1 与定理 5.1.5 的结论对以 $A$ 为系数的范畴仍然成立，即有范畴的等价
\begin{center}
\begin{tikzcd}
\mathrm{VB}^{\Gamma}_{X_{F_\infty,A}} \arrow[r] \arrow[d, "\sim"] & \mathrm{VB}^{\varphi,\Gamma}_{Y_{F_\infty,A}} \arrow[d, "\sim"] & \Mod^{\varphi,\Gamma}_{\mathcal{R}_{F,A}} \arrow[l] \\
\mathrm{VB}^{\mathrm{Gal}_F}_{\overline{X}_{F,A}} \arrow[r, "\sim"] & \mathrm{VB}^{\varphi,\mathrm{Gal}_F}_{\overline{Y}_{F,A}}
\end{tikzcd}
\end{center}
这些等价由与无系数情形类似的函子诱导。此外，给定仿射代数的态射 $A\to B$，有显然的基变换函子
\[
(-)\widehat{\otimes}_AB:\mathrm{VB}^{\Gamma}_{X_{F_\infty,A}}\longrightarrow\mathrm{VB}^{\Gamma}_{X_{F_\infty,B}}
\]
等等，且上述范畴的等价显然与从 $A$ 到 $B$ 的基变换相容。

\begin{rem}{5.1.6}\label{rem:a4-5-1-6}
还值得指出，用 Beauville--Laszlo 粘合引理对这些等价范畴所作的描述可以推广到族的情形。这使我们能用沿 $\{\infty\}\times\Sp A$ 的完备化上的一族不变格，去修改 $\overline{X}_{F,A}$ 上给定的一族等变向量丛。
\end{rem}

\subsubsection*{5.1.7. 刚性解析 Artin 叠}
我们在刚性解析空间范畴上的 Artin 叠框架中引入基本概念，为定义刚性解析 $(\varphi,\Gamma)$-模的叠作准备。记 $\mathrm{Rig}_L$ 为固定基域 $L$ 上的刚性解析空间范畴，其中 $L$ 是 $\Q_p$ 的有限扩张。我们将给 $\mathrm{Rig}_L$ 配上 Tate-fpqc 拓扑 [CT09, 2.1]。该拓扑的覆盖由通常的（容许）Tate 覆盖，以及仿射代数忠实平坦映射 $B\to A$ 所对应的刚性空间态射 $\Sp(A)\to\Sp(B)$ 生成。关于该拓扑，一切可表示函子都是层，且凝聚层满足下降 [Con06, Theorem 4.2.8]。下文所说的 $\mathrm{Rig}_L$ 上的叠，是指 $\mathrm{Rig}_L$ 上纤维于群胚、且对 Tate-fpqc 拓扑满足下降的范畴。与概形上 Artin 叠的情形一样，我们必须从代数空间的类比物的定义开始。

\begin{defn}{5.1.8}\label{defn:a4-5-1-8}
拟解析空间（quasi-analytic space）是 $\mathrm{Rig}_L$ 上的层 $\mathcal{F}$，使得对角线 $\mathcal{F}\to\mathcal{F}\times\mathcal{F}$ 可表示，且存在从可表示层 $U$ 到 $\mathcal{F}$ 的 étale 满射 $U\to\mathcal{F}$。
\end{defn}

\ysrc{206}
\begin{rem}{5.1.9}\label{rem:a4-5-1-9}
\begin{enumerate}[label=(\roman*),leftmargin=2.4em,itemsep=0.2em]
\item 事实上，下文并不需要这种一般性：由 Conrad 与 Temkin 的一个结果 [CT09, Theorem 1.2.2]，每个分离的拟解析空间都可由刚性解析空间表示。下文只会遇到分离的拟解析空间。
\item Evan Warner 的博士论文 [War17] 也发展了一套 adic 空间上 Artin 叠的理论。它与我们框架的主要差别在于：Warner 处理全部强诺特 adic 空间（并使用 étale 拓扑），而我们只允许刚性解析空间（并使用 Tate-fpqc 拓扑）；此外 Warner 使用了不同的术语。我们主要限于刚性解析空间，是因为至少目前，带系数的 $(\varphi,\Gamma)$-模理论 [KPX14]，特别是关于其上同调的有限性结果，只对经典刚性解析几何中仿射代数的系数有定义。
\end{enumerate}
\end{rem}

\begin{defn}{5.1.10}\label{defn:a4-5-1-10}
刚性解析 Artin 叠是 $\mathrm{Rig}_L$ 上的叠 $\mathfrak{X}$，使得对角线 $\Delta_{\mathfrak{X}}:\mathfrak{X}\to\mathfrak{X}\times_L\mathfrak{X}$ 可由拟解析空间表示，且存在刚性解析空间 $U$ 与光滑满射 $U\to\mathfrak{X}$。
\end{defn}

给定刚性解析 Artin 叠 $\mathfrak{X}$，可按通常方式定义其凝聚层范畴
\[
\Coh(\mathfrak{X})=\varprojlim_{\Sp(A)\to\mathfrak{X}}\Coh(\Sp(A)),
\]
其中极限取遍仿射刚性空间 $\Sp(A)$ 到 $\mathfrak{X}$ 的一切映射。按通常方式，给定从（仿射）刚性解析空间 $U$ 到 $\mathfrak{X}$ 的光滑满射 $U\to\mathfrak{X}$，该范畴可计算为图
\[
\Coh(U)\rightrightarrows\Coh(U\times_{\mathfrak{X}}U)\rightrightarrows\rightrightarrows\cdots
\]
的极限。类似地，$\mathfrak{X}$ 上凝聚层的稳定 $\infty$-范畴可定义为 $\infty$-范畴的同伦极限
\[
\mathbf{D}_{\mathrm{coh}}(\mathfrak{X})=\varprojlim_{\Sp(A)\to\mathfrak{X}}\mathbf{D}_{\mathrm{coh}}(\Sp(A)),
\]
而对给定的覆盖，它可计算为 $\infty$-范畴的图
\[
\mathbf{D}_{\mathrm{coh}}(U)\rightrightarrows\mathbf{D}_{\mathrm{coh}}(U\times_{\mathfrak{X}}U)\rightrightarrows\rightrightarrows\cdots
\]
的同伦极限。

对刚性解析空间的态射 $f:X\to Y$，更一般地对刚性解析 Artin 叠之间的态射，上述定义蕴含存在典范的导出拉回 $Lf^*:\mathbf{D}_{\mathrm{coh}}(Y)\to\mathbf{D}_{\mathrm{coh}}(X)$。若 $f$ 是固有映射，还有（导出的）推前 $Rf_*:\mathbf{D}_{\mathrm{coh}}(X)\to\mathbf{D}_{\mathrm{coh}}(Y)$，它是 $Lf^*$ 的右伴随。更确切地说，对一般的 $f$，推前 $Rf_*$ 可定义在 $\mathbf{D}_{\mathrm{coh}}(X)$ 中在 $Y$ 上具有固有支撑的层所构成的子范畴上。对我们所考虑的刚性解析 Artin 叠，总存在对偶复形 $\omega_{\mathfrak{X}}$，而内蕴对偶
\[
\mathbf{D}_{\mathfrak{X}}(-)=R\mathcal{H}om_{\mathcal{O}_{\mathfrak{X}}}(-,\omega_{\mathfrak{X}})
\]
使我们能把
\[
f^!=\mathbf{D}_{\mathfrak{X}}\circ Lf^*\circ\mathbf{D}_{\mathfrak{Y}}:\mathbf{D}_{\mathrm{coh}}(\mathfrak{Y})\to\mathbf{D}_{\mathrm{coh}}(\mathfrak{X})
\]
当作上 shriek 拉回（upper shriek pullback）$f^!$ 的权宜定义。

\ysrc{207}
\begin{rem}{5.1.11}\label{rem:a4-5-1-11}
上述定义有些局限，因为我们只定义了凝聚层范畴，而未定义拟凝聚层范畴。这主要是因为刚性解析空间上的拟凝聚层（出于拓扑原因）较为微妙。在 Clausen--Scholze 的凝聚态环与模的世界里，可以发展（导出的）拟凝聚层范畴的理论，并配以完整的六函子形式，这比上面关于 $f^!$ 的权宜定义更为统一。
\end{rem}

我们提一下刚性解析 Artin 叠的两个来源。给定 $L$-概形范畴 $\mathrm{Sch}_L$ 上的 Artin 叠 $\mathbf{X}$，且 $\mathbf{X}$ 在 $L$ 上（局部）有限型（即 $\mathbf{X}$ 可由一个在 $L$ 上（局部）有限型的概形 $\mathbf{U}$ 覆盖），其解析化可定义为
\[
\mathbf{X}^{\mathrm{an}}(\Sp(A))=\mathbf{X}(\Spec(A))
\]
的层化。容易看出，在 $\mathbf{X}=\mathbf{A}^n$ 的情形，从而在（有限型）仿射概形的情形，这就是正确的定义。特别地：对表现 $\mathbf{V}=\mathbf{U}\times_{\mathbf{X}}\mathbf{U}\rightrightarrows\mathbf{U}\to\mathbf{X}$，$\mathbf{X}$ 的解析化可定义为 $\mathbf{V}^{\mathrm{an}}\rightrightarrows\mathbf{U}^{\mathrm{an}}$ 的叠商。给定 $\mathbf{X}$，在（导出）凝聚层范畴上有显然的解析化函子，它与 $Lf^*$、$f^!$、$Rf_*$（只要它们有定义）交换。

另一方面，若 $\mathcal{X}$ 是 $\Spf(\mathcal{O}_L)$ 上 [Eme, 5.] 意义下的形式代数叠，则可定义其一般纤维为
\begin{equation}
\mathcal{X}^{\mathrm{rig}}_\eta(\Sp(A))=\varinjlim_{Y}\mathcal{X}(Y),\tag{5.1.12}
\end{equation}
其中极限取遍 $\Sp(A)$ 的一切（Raynaud 意义下的）形式模型 $Y$。

\begin{prop}{5.1.13}\label{prop:a4-5-1-13}
设 $\mathcal{X}$ 是拓扑有限型、且对角线拓扑有限型的形式代数叠。则 $\mathcal{X}^{\mathrm{rig}}_\eta$ 是刚性解析 Artin 叠。
\end{prop}

这里，称形式代数叠 $\mathcal{X}$ 拓扑有限型，若它容许从某个形式概形的光滑满射，而该形式概形局部形如 $\Spf A$，其中 $A$ 是拓扑有限型的完备 $\Z_p$-代数（对角线拓扑有限型的定义类似）。证明命题 5.1.13 的想法是归约到仿射情形，此时 $\mathcal{X}$ 可写成形式光滑映射
\[
\Spf B\rightrightarrows\Spf A
\]
的余等化子，其中 $A,B$ 是拓扑有限型的完备 $\Z_p$-代数。于是 $\mathcal{X}$ 的刚性解析一般纤维可等同于
\[
(\Spf B)^{\mathrm{rig}}_\eta\rightrightarrows(\Spf A)^{\mathrm{rig}}_\eta
\]
的余等化子，其中 $(-)^{\mathrm{rig}}_\eta$ 表示形式概形在 Berthelot 意义下的一般纤维。

\subsubsection*{5.1.14. 刚性解析 $(\varphi,\Gamma)$-模的叠}
现在我们回到本小节的主要目标，定义 $(\varphi,\Gamma)$-模的叠。

\begin{defn}{5.1.15}\label{defn:a4-5-1-15}
对 $d\ge1$，$\mathrm{Rig}_L$ 上纤维于群胚的范畴 $\mathfrak{X}_{\mathrm{GL}_d}=\mathfrak{X}_d$ 定义如下：它把仿射空间 $\Sp A$ 对应到 $\mathcal{R}_{F,A}$ 上秩为 $d$ 的 $(\varphi,\Gamma)$-模构成的群胚。
\end{defn}

\ysrc{208}
利用 $\mathcal{R}_{F,A}$ 上 $(\varphi,\Gamma)$-模在 $\Sp A$ 上的圆环上有模型这一事实，刚性解析空间上向量丛的下降给出：

\begin{prop}{5.1.16}\label{prop:a4-5-1-16}
纤维于群胚的范畴 $\mathfrak{X}_d$ 是叠。
\end{prop}

\begin{rem}{5.1.17}\label{rem:a4-5-1-17}
关于 $\mathfrak{X}_d$ 与第 4 节的叠 $\mathcal{X}_d$ 之间预期的关系，见 (5.1.31) 及周围的讨论。
\end{rem}

\begin{rem}{5.1.18}\label{rem:a4-5-1-18}
由上述讨论，叠 $\mathfrak{X}_d$ 可等价地描述为 Fargues--Fontaine 曲线上 $\mathrm{Gal}_F$-等变向量丛的叠；下文我们会自由选用更适合目的的哪一种描述。于是它与代数曲线上（向量丛的）叠的经典理论有明显的相似之处。我们列出最重要的相似与差异：
\begin{enumerate}[label=(\roman*),leftmargin=2.6em,itemsep=0.2em]
\item 有上同调理论
\[
R\Gamma_{\mathrm{Gal}_F}(X,-):=R\Gamma(\mathrm{Gal}_F,R\Gamma(\Sp A\times\overline{X}_F,-))
\]
用于 $\Sp A\times\overline{X}_F$ 上的 $\mathrm{Gal}_F$-等变向量丛。就 $(\varphi,\Gamma)$-模而言，该上同调理论与由 Herr 复形计算的通常 $(\varphi,\Gamma)$ 上同调一致 [KPX14, Def. 2.2.3]。更确切地说，给定 $\mathcal{R}_{F,A}$ 上的 $(\varphi,\Gamma)$-模 $D$，相应的 Herr 复形为
\[
R\Gamma_{\varphi,\Gamma}(D)=\bigl[D^\Delta\xrightarrow{(\varphi-1,\gamma-1)}D^\Delta\oplus D^\Delta\xrightarrow{(\gamma-1)\oplus(1-\varphi)}D^\Delta\bigr],
\]
其中 $\Delta\subset\Gamma$ 是挠子群，而 $\gamma\in\Gamma/\Delta$ 是拓扑生成元。该复形（拟同构于）一个集中在次数 $0,1,2$ 的 $A$-模完美复形，见 [KPX14, Theorem 4.4.5]。
\item 作为 (i) 的推论，Grothendieck 关于凝聚层的上同调与基变换定理，在 Fargues--Fontaine 曲线上等变向量丛的框架中（相应地，在 $(\varphi,\Gamma)$-模的框架中）有类比物。更确切地说，给定 $\mathcal{R}_{F,A}$ 上的 $(\varphi,\Gamma)$-模 $D$ 与点 $x\in\Sp A$，基变换映射
\[
\mathrm{bc}^i_x:H^i_{\varphi,\Gamma}(D)\otimes_Ak(x)\longrightarrow H^i_{\varphi,\Gamma}(D\otimes_{\mathcal{R}_{F,A}}\mathcal{R}_{F,k(x)})
\]
是满射，当且仅当它是同构；且若 $\mathrm{bc}^i_x$ 是满射，则 $H^i_{\varphi,\Gamma}(D)$ 在 $x$ 的某个邻域内是向量丛，当且仅当 $\mathrm{bc}^{i-1}_x$ 是满射。
\item 没有丰沛线丛（ample line bundle）的类比物：对给定的等变向量丛 $V$，Euler 示性数公式（见前引文献）
\[
\sum_{i\ge0}(-1)^i\dim H^i_{\mathrm{Gal}_F}(\overline{X}_F,V)=-\mathrm{rk}\,V[F:\Q_p]
\]
蕴含 $H^1$ 总是不消失（当然除非 $V=0$）。特别地，我们无法通过扭掉高阶上同调。（若不考虑 $\mathrm{Gal}_F$-等变性，$\overline{X}_F$ 上有丰沛线丛的类比物。这意味着我们至少可以扭向量丛 $V$，使其 $H^2$ 消失。）
\item \ysrc{209}局部 Tate 对偶 [KPX14, Theorem 4.4.5 (3)] 断言：给定 $\mathcal{R}_{F,A}$ 上的 $(\varphi,\Gamma)$-模，有典范同构
\[
R\Gamma_{\varphi,\Gamma}(D)\xrightarrow{\ \sim\ }R\Hom_A\bigl(R\Gamma_{\varphi,\Gamma}(D^\vee(\varepsilon)),A\bigr)[-2].
\]
这里 $D^\vee=\Hom_{\mathcal{R}_{F,A}}(D,\mathcal{R}_{F,A})$，$\varepsilon$ 是分圆特征标。
\item 有一套 Hodge--Tate--Sen 权理论，我们在下文 §5.1.33 中解释。它在 §5.3 关于抛物子群 $P\subseteq\mathrm{GL}_d$ 的 $P$-结构的讨论中引入实质性的差异。
\end{enumerate}
\end{rem}

\begin{conj}{5.1.19}\label{conj:a4-5-1-19}
叠 $\mathfrak{X}_d$ 是维数为 $d^2[F:\Q_p]=\dim\operatorname{Res}_{F/\Q_p}\mathrm{GL}_d$ 的刚性解析 Artin 叠，并且是正规的、局部完全交（local complete intersection）的。
\end{conj}

下面描述朝这个猜想方向的一些部分结果。

\begin{thm}{5.1.20}\label{thm:a4-5-1-20}
\begin{enumerate}[label=(\roman*),leftmargin=2.4em,itemsep=0.2em]
\item $\mathfrak{X}_d$ 的对角线 $\Delta_{\mathfrak{X}_d}:\mathfrak{X}_d\to\mathfrak{X}_d\times\mathfrak{X}_d$ 可由刚性空间表示。
\item 设 $L'$ 是 $L$ 的有限扩张，$x\in\mathfrak{X}_d(L')$ 是一个 $L'$-值点。则存在 $x$ 在 $\mathfrak{X}_d$ 中的开邻域 $\mathcal{U}_x$，它是刚性解析 Artin 叠。此外，若 $d=2$，则 $\mathcal{U}_x$ 具有预期的维数，并且正规、是局部完全交。
\end{enumerate}
\end{thm}

\begin{proof}[证明梗概]
(i) 这基本上由 $R\Gamma_{\varphi,\Gamma}(D)$ 总可由一个完美复形泛地计算这一事实推出。更确切地说，给定仿射刚性解析空间 $\Sp A$ 上两族 $(\varphi,\Gamma)$-模 $D_1,D_2$，要证明函子
\[
\underline{\mathrm{Isom}}(D_1,D_2):\Sp B\longmapsto\underline{\mathrm{Isom}}_{\varphi,\Gamma}(D_1\widehat{\otimes}_AB,D_2\widehat{\otimes}_AB)
\]
在 $\Sp A$ 上的仿射空间范畴中可由刚性解析空间表示。

考察最高外幂上的诱导映射可知，$\underline{\mathrm{Isom}}(D_1,D_2)$ 是
\[
\underline{\Hom}(D_1,D_2)=\underline{H}^0_{\varphi,\Gamma}(D_1^\vee\otimes D_2):\Sp B\longmapsto\underline{\Hom}_{\varphi,\Gamma}(D_1\widehat{\otimes}_AB,D_2\widehat{\otimes}_AB)
\]
\[
{}=\underline{H}^0_{\varphi,\Gamma}\bigl((D_1^\vee\otimes D_2)\widehat{\otimes}_AB\bigr)
\]
的开子函子。记 $D=D_1^\vee\otimes D_2$。则存在集中在次数 $[0,2]$ 的向量丛复形 $\mathcal{E}^\bullet$，以及同构 $\underline{H}^0_{\varphi,\Gamma}(D)\cong\ker(\mathcal{E}^0\to\mathcal{E}^1)$，其中右端函子把 $\Sp B$ 送到 $\ker(\mathcal{E}^0\otimes_AB\to\mathcal{E}^1\otimes_AB)$。由于 $\mathcal{E}^0$ 与 $\mathcal{E}^1$ 是向量丛，该函子可由 $\Sp A$ 上相应几何向量丛之间诱导映射的零截面的原像表示。

(ii) 我们仅对 $d=2$ 的情形勾勒证明。显然可归约到 $L'=L$。设 $D$ 是 $\mathcal{R}_{F,L}$ 上秩为 $2$ 的 $(\varphi,\Gamma)$-模。为构造一个卡（chart），即从某个刚性解析空间到 $D$ 的邻域的光滑满射，区分两种可能情形：
\begin{enumerate}[label=(\alph*),leftmargin=2.2em,itemsep=0.15em]
\item 存在特征 $\delta:F^\times\to L^\times$，使得 $D'=D\otimes_{\mathcal{R}_{F,L}}(\delta)$ 是 étale 的。
\item 不存在这样的特征。
\end{enumerate}
在情形 (a) 中，可归约到伽罗瓦表示的情形，并利用（加框架）伽罗瓦形变环的一般纤维来构造一个卡\ysrc{210}，该卡在 $D$ 处局部有效。在情形 (b) 中，我们同样想归约到伽罗瓦表示的情形：在用某个特征扭 $D$ 之后，构造一个扩张
\begin{equation}
0\to D\to D'\to\mathcal{R}_{F,L}(\delta)\to0,\tag{5.1.21}
\end{equation}
使得 $D'$ 是附于某个 $3$ 维伽罗瓦表示 $\rho'$ 的 $(\varphi,\Gamma)$-模。于是 $\rho'$ 定义了某个剩余 $3$ 维伽罗瓦表示 $\overline{\rho}'$ 的（加框架）形变空间的一般纤维 $(\Spf R_{\overline{\rho}'})^{\mathrm{rig}}$ 中的一点，而（在 $\rho'$ 局部附近）空间 $(\Spf R_{\overline{\rho}'})^{\mathrm{rig}}$ 含有一个闭子空间 $Y$，在其上可以形变扩张 (5.1.21)。把 $Y$ 上的泛伽罗瓦表示映到此扩张中的子对象，便定义了所要的到 $\mathfrak{X}_2$ 的态射；我们将证明它在 $\rho'$ 处光滑。下面更详细地解释情形 (a) 与 (b) 的论证：

情形 (a)：由于 $D'$ 是 étale 的，存在伽罗瓦表示 $\rho':\mathrm{Gal}_F\to\mathrm{GL}_2(L)$ 使 $D_{\mathrm{rig}}(\rho')=D'$。必要时用 $\mathrm{GL}_2(L)$ 中某个元素共轭 $\rho'$，可假定 $\rho'$ 取值于 $\mathrm{GL}_2(\mathcal{O})$（即我们在伽罗瓦表示中选取一个稳定格）。记 $\overline{\rho}':\mathrm{Gal}_F\to\mathrm{GL}_2(k)$ 为 $\rho'$ 的约化，$R_{\overline{\rho}'}$ 为 $\overline{\rho}'$ 的泛加框架形变环。则 $\rho'$ 定义了刚性解析一般纤维 $(\Spf R_{\overline{\rho}'})^{\mathrm{rig}}$ 中一点，剩下只需证明态射
\[
(\Spf R_{\overline{\rho}'})^{\mathrm{rig}}\longrightarrow\mathfrak{X}_2,\qquad
\rho\longmapsto D_{\mathrm{rig}}(\rho)\otimes\mathcal{R}(\delta^{-1})
\]
（它把 $\rho'$ 送到 $D$）在 $\rho'$ 的邻域内光滑。这由无穷小提升判别法，连同 $(\Spf R_{\overline{\rho}'})^{\mathrm{rig}}$ 在 $\rho'$ 处的完备局部环前表示 $\rho':\mathrm{Gal}_F\to\mathrm{GL}_2(L)$ 的泛加框架形变函子这一事实推出，见 [Kis09c, Lemma 2.3.3, Prop. 2.3.5]。

情形 (b)：再用某个秩 $1$ 对象作扭之后，可假定 $D$ 的一切子对象的斜率都非负。由下文引理 5.1.23，存在特征 $\delta:F^\times\to L^\times$ 与非分裂扩张
\[
0\longrightarrow D\longrightarrow D'\longrightarrow\mathcal{R}_{F,L}(\delta)\longrightarrow0,
\]
使得 $D'$ 是 étale 的，且
\[
\Ext^0_{\varphi,\Gamma}(\mathcal{R}_{F,L}(\delta),D)=\Ext^2_{\varphi,\Gamma}(\mathcal{R}_{F,L}(\delta),D)=0,
\]
并有 $\dim\Hom_{\varphi,\Gamma}(D',\mathcal{R}_{F,L}(\delta))=1$。我们再次在附于 $D'$ 的伽罗瓦表示中选取一个 $\mathcal{O}$-格 $\rho'$，并把 $\rho'$ 看作 $(\Spf R_{\overline{\rho}'})^{\mathrm{rig}}$ 中一点，后者是 $\rho'$ 的约化 $\overline{\rho}'$ 的泛加框架形变环的刚性解析一般纤维。

如下构造含 $\rho'$ 的局部闭子集 $Y\subseteq(\Spf R_{\overline{\rho}'})^{\mathrm{rig}}$ 以及光滑态射 $Y\to\mathfrak{X}_2$：记 $\tilde{D}$ 为 $(\Spf R_{\overline{\rho}'})^{\mathrm{rig}}$ 上的泛 $(\varphi,\Gamma)$-模，$Y_1$ 为 $H^2_{\varphi,\Gamma}(\tilde{D}(\delta^{-1}\varepsilon))$ 的概形论支撑。由于
\[
H^2_{\varphi,\Gamma}(\tilde{D}(\delta^{-1}\varepsilon))\otimes k(\rho')=H^2_{\varphi,\Gamma}(D'(\delta^{-1}\varepsilon))=\bigl(\Hom_{\varphi,\Gamma}(D',\mathcal{R}_{F,L}(\delta))\bigr)^\vee
\]
的维数为 $1$，可找到 $\rho'$ 的邻域 $Y_2\subseteq Y_1$ 以及 $\Gamma(Y_2,H^2_{\varphi,\Gamma}(\tilde{D}(\delta^{-1}\omega)))$ 中元素 $\tilde{f}$，它是某个选定基向量
\[
f^\vee\in\Hom_{\varphi,\Gamma}(D',\mathcal{R}_{F,L}(\delta))
\]
的对偶基的提升\tn{底本此处作 $\delta^{-1}\omega$，而本段上下文与 (5.1.22) 均作 $\delta^{-1}\varepsilon$；$\omega$ 与 $\varepsilon$ 都是分圆特征标的常用记号，疑为笔误。}。必要时用 $\rho'$ 的更小开邻域替换 $Y_2$，可假定 $H^2_{\varphi,\Gamma}(\tilde{D}(\delta^{-1}\varepsilon))|_{Y_2}$ 是以 $\tilde{f}$ 为基的自由秩 $1$ 模。再作局部化，\ysrc{211}便找到 $\rho'$ 的邻域 $Y\subseteq Y_2$ 以及满射
\[
\tilde{f}^\vee:\tilde{D}|_Y\longrightarrow\mathcal{R}_{F,Y}(\delta),
\]
它在 $\rho'$ 处特殊化为所选满射 $f^\vee:D'\to\mathcal{R}_{F,L}(\delta)$；此外还可假定对一切 $y\in Y$ 都有
\begin{equation}
\Ext^0\bigl(\mathcal{R}_{F,k(y)}(\delta),\tilde{D}\otimes k(y)\bigr)=\Ext^2\bigl(\mathcal{R}_{F,k(y)}(\delta),\tilde{D}\otimes k(y)\bigr)=0.\tag{5.1.22}
\end{equation}
我们断言，由 $\tilde{D}|_Y\mapsto\ker\tilde{f}^\vee$ 定义的映射 $Y\to\mathfrak{X}_2$ 光滑。同样只需在点 $y\in Y$ 的完备局部环上验证。$y$ 处的完备局部环前表示如下函子：$\tilde{D}_y=\tilde{D}\otimes k(y)$ 的形变，连同映射 $f^\vee_y:\tilde{D}_y\to\mathcal{R}_{F,k(y)}(\delta)$ 的形变。由我们对 $Y$ 的选取，该函子可等同于 $D_y=\ker f^\vee_y$ 的形变，连同扩张
\[
0\longrightarrow D_y\longrightarrow\tilde{D}_y\longrightarrow\mathcal{R}_{F,L}(\delta)\longrightarrow0
\]
的形变所构成的函子。由 Ext-群 (5.1.22) 的消没，该函子在 $D_y$ 的形变函子上形式光滑。

剩下要证关于维数的断言与关于局部结构的断言，这就是下文命题 5.1.27。
\end{proof}

\begin{lem}{5.1.23}\label{lem:a4-5-1-23}
设 $D$ 是 $\mathcal{R}_{F,L}$ 上秩为 $2$ 的 $(\varphi,\Gamma)$-模，假定 $D$ 的一切子对象的斜率都非负，且对任一特征 $\delta:F^\times\to L^\times$ 都有 $D\not\cong\mathcal{R}_{F,L}(\delta)^{\oplus2}$。设 $\eta:F^\times\to L^\times$ 是连续特征，满足
\begin{itemize}[leftmargin=1.8em,itemsep=0.15em]
\item 对一切嵌入 $\tau:F\hookrightarrow L$ 都有 $(\mathrm{wt}_\tau(\eta)+\Z)\cap\mathrm{wt}_\tau(D)=\varnothing$，其中权的定义见下文 §5.1.33；
\item $\eta$ 的斜率为 $\mathrm{slope}(\eta)=-\mathrm{slope}(D)$。
\end{itemize}
则存在非分裂扩张
\begin{equation}
0\longrightarrow D\longrightarrow D'\longrightarrow\mathcal{R}_{F,L}(\eta)\longrightarrow0,\tag{5.1.24}
\end{equation}
使得 $D'$ 是 étale 的，并且
\[
\Ext^0_{\varphi,\Gamma}(\mathcal{R}_{F,L}(\eta),D)=\Ext^2_{\varphi,\Gamma}(\mathcal{R}_{F,L}(\eta),D)=0,\qquad
\dim_L\Hom_{\varphi,\Gamma}(D',\mathcal{R}_{F,L}(\eta))=1.
\]
\end{lem}

\begin{rem}{5.1.25}\label{rem:a4-5-1-25}
给定引理中的 $D$，显然总存在满足引理假设的特征 $\eta$。
\end{rem}

\begin{proof}
由于 $D(\eta^{-1})$ 与 $D^\vee(\eta\omega)$ 具有非整的 Hodge--Tate--Sen 权，Ext-群的消没是下述等同的直接推论：
\[
\Ext^0_{\varphi,\Gamma}(\mathcal{R}_{F,L}(\eta),D)=H^0_{\varphi,\Gamma}(D(\eta^{-1})),\qquad
\Ext^2_{\varphi,\Gamma}(\mathcal{R}_{F,L}(\eta),D)=H^0_{\varphi,\Gamma}(D^\vee(\eta\omega))^\vee.
\]
任取一个非分裂扩张 (5.1.24)。关于 Hodge--Tate 权的断言也蕴含关于 $\dim\Hom_{\varphi,\Gamma}(D',\mathcal{R}_{F,L}(\eta))$ 的断言，因此只需说明可选 $D'$ 为 étale。由关于 $\eta$ 斜率的条件，自动有 $\mathrm{slope}(D')=0$；再由 Kedlaya 的斜率滤链定理 [Ked08]，只需说明对一切（饱和）子对象 $D''\subseteq D$ 都有 $\mathrm{slope}(D'')\ge0$。按假定，当 $D''\subseteq D$ 时这是对的。因此只需说明可选 $D'$，使 $D'$ 的一切饱和 $(\varphi,\Gamma)$-稳定子对象都含于 $D$。反设存在 $D''$，使映射 $D''\to\mathcal{R}_{F,L}(\eta)$ 非零。区分 $\mathrm{rk}\,D''=1$ 与 $\mathrm{rk}\,D''=2$ 两种情形。
\end{proof}

\ysrc{212}
若 $\mathrm{rk}\,D''=1$，则 $D''\to\mathcal{R}_{F,L}(\eta)$ 在反转 $t$ 之后是同构，这蕴含扩张 (5.1.24) 在反转 $t$ 之后分裂。然而假定 $(\mathrm{wt}_\tau(\eta)+\Z)\cap\mathrm{wt}_\tau(D)=\varnothing$ 蕴含典范映射
\[
\Ext^1_{\varphi,\Gamma}(\mathcal{R}_{F,L},D)\longrightarrow\Ext^1_{\varphi,\Gamma}(\mathcal{R}_{F,L}[1/t],D[1/t])
\]
是同构，这与我们选取 (5.1.24) 为非分裂相矛盾。

若 $\mathrm{rk}\,D''=2$，则 $E=D''\cap D\subseteq D$ 是秩 $1$ 的 $(\varphi,\Gamma)$-子模。按假定，$D$ 只有有限多个秩 $1$ 的饱和子对象。每个这样的子对象 $E$ 给出短正合列
\begin{equation}
0\longrightarrow\mathcal{R}_{F,L}(\delta_1)\longrightarrow D\longrightarrow\mathcal{R}_{F,L}(\delta_2)\longrightarrow0,\tag{5.1.26}
\end{equation}
而我们关于 Hodge--Tate--Sen 权的假定蕴含某些 $\Ext^0$ 与 $\Ext^2$ 群的消没，从而给出短正合列
\[
0\to\Ext^1_{\varphi,\Gamma}(\mathcal{R}_{F,L}(\eta),\mathcal{R}_{F,L}(\delta_1))\to\Ext^1_{\varphi,\Gamma}(\mathcal{R}_{F,L}(\eta),D)\to
\]
\[
{}\to\Ext^1_{\varphi,\Gamma}(\mathcal{R}_{F,L}(\eta),\mathcal{R}_{F,L}(\delta_2))\to0.
\]
此外，与秩 $1$ 的情形一样，条件 $(\mathrm{wt}_\tau(\eta)+\Z)\cap\mathrm{wt}_\tau(D)=\varnothing$ 蕴含反转 $t$ 在该正合列与其对应正合列
\[
0\to\Ext^1_{\varphi,\Gamma}(\mathcal{R}_{F,L}(\eta)[1/t],\mathcal{R}_{F,L}(\delta_1)[1/t])\to\Ext^1_{\varphi,\Gamma}(\mathcal{R}_{F,L}(\eta)[1/t],D[1/t])\to
\]
\[
{}\to\Ext^1_{\varphi,\Gamma}(\mathcal{R}_{F,L}(\eta)[1/t],\mathcal{R}_{F,L}(\delta_2)[1/t])\to0
\]
之间诱导同构。由于 $D$ 只有有限多个子对象，我们可以选取扩张
\[
D'\in\Ext^1_{\varphi,\Gamma}(\mathcal{R}_{F,L}(\eta),D),
\]
使得对把 $D$ 写成扩张 (5.1.26) 的一切可能方式，$D'$ 都不在 $\Ext^1_{\varphi,\Gamma}(\mathcal{R}_{F,L}(\eta),\mathcal{R}_{F,L}(\delta_1))$ 的像中。此时，若 $D''\subseteq D'$ 是饱和子对象，满足 $D''\cap D=E=\mathcal{R}_{F,L}(\delta_1)$，且 $D''\to\mathcal{R}_{F,L}(\eta)$ 非零，则
\[
\mathcal{R}_{F,L}(\eta)[1/t]\cong(D''/E)[1/t]\subseteq(D'/E)[1/t]
\]
蕴含 $D'$ 在 $\Ext^1_{\varphi,\Gamma}(\mathcal{R}_{F,L}(\eta)[1/t],\mathcal{R}_{F,L}(\delta_2)[1/t])$ 中的像 $(D'/E)[1/t]$ 消失。这与上面关于 $D'$ 的选取矛盾。

局部卡的存在蕴含 $\mathfrak{X}_d$ 在刚性解析点处 versal 环的存在。

\begin{prop}{5.1.27}\label{prop:a4-5-1-27}
设 $d=2$，$x\in\mathfrak{X}_d$，$R$ 是 $x$ 处的 versal 环，使得 $\Spf R\to\mathfrak{X}_d$ 是相对维数为 $m$ 的形式光滑映射。则 $R$ 是维数为 $m+d^2[F:\Q_p]$ 的局部完全交，并且正规。
\end{prop}

\begin{proof}[证明梗概]
由于任意两个 versal 环光滑等价，可用任一 versal 环来证明该断言。我们沿用定理 5.1.20 证明中的记号。若 $x$ 处于该证明的情形 (a)，则 $R$ 可取为 $(\Spf R_{\overline{\rho}'})^{\mathrm{rig}}$ 在点 $\rho'$ 处的完备局部环，结论由 [BIP23a] 推出。在情形 (b)，即 $x$ 不是某个 étale $(\varphi,\Gamma)$-模的扭时，标准的形变理论论证给出 versal 环 $R$ 的一个表现，使得 $\Spf R\to\mathfrak{X}_d$ 是维数为 $m$ 的形式光滑映射，且 $R$ 的每个不可约分支的维数至少为 $m+[F:\Q_p]d^2$；\ysrc{213}为说明 $R$ 是局部完全交，需证对一切不可约分支等号成立。由局部卡的存在性（即定理 5.1.20 (2)），可把 $R$ 看作某点 $y\in Y=\Sp A$ 的完备局部环，其中 $f:Y=\Sp A\to\mathfrak{X}_d$ 是相对维数为 $m$ 的光滑映射。限制到 $Y$ 的开子集之后，可假定 $Y$ 的一切不可约分支都含 $y$，于是只需说明 $Y$ 的一切不可约分支维数至多为 $m+[F:\Q_p]d^2$。设 $Z$ 是这样的不可约分支，则在 $Z$ 的一般点处，即在某个 Zariski 开子集 $U_Z\subset Z$ 上，限制 $f_Z:Z\to\mathfrak{X}_d$ 是相对维数为 $m$ 的形式光滑映射。若 $U_Z$ 含有点 $z$，使得 $(\varphi,\Gamma)$-模 $f_Z(z)$ 在扭的意义下是 étale 的，则由情形 (a) 的讨论，$U_Z$ 在 $z$ 处的完备局部环维数为 $m+[F:\Q_p]d^2$，从而 $Z$ 的维数为 $m+[F:\Q_p]d^2$。由于秩 $2$ 的 $(\varphi,\Gamma)$-模若在扭的意义下不是 étale 的则必可约（此处用到 Kedlaya 的斜率滤链定理 [Ked08]，它断言「在扭的意义下 étale」与「半稳定」是一回事），这意味着我们必须排除如下可能：（在不可约分支 $Z$ 上的一般点处）一切点 $z$ 都具有使 $(\varphi,\Gamma)$-模 $f(z)$ 可约的性质。

由于 $Z$ 的维数至少为 $m+[F:\Q_p]d^2$，且它（在一般点处）是 $\mathfrak{X}_d$ 上相对维数为 $m$ 的光滑映射，我们可以通过证明秩 $2$ 的可约 $(\varphi,\Gamma)$-模构成的空间维数太小来排除这一可能。我们在形变空间的层面上勾勒论证\footnote{严格说来，需要证明的略强一些，即：刚性解析空间 $Z$ 中使 $(\varphi,\Gamma)$-模逐点可约的子集，是维数不超过 $m+[F:\Q_p]\frac{d(d+1)}{2}<m+[F:\Q_p]d^2$ 的 Zariski 闭子空间的至多可数并。这些陈述可与下文定理 5.3.12 及注 5.3.15 比较。}。给定扩张 (5.1.26)，考虑该短正合列的形变构成的群胚 $\mathfrak{D}$，记 $S$ 为 $\mathfrak{D}$ 的 versal 环，使得 $\Spf S\to\mathfrak{D}$ 是维数为 $m$ 的形式光滑映射。于是只需说明 $S$ 的一切不可约分支维数严格小于 $m+[F:\Q_p]d^2$。在 $d=2$ 的情形，这可由直接计算容易地验证（也可与下文 5.3 中关于 $B$-丛的叠的讨论比较）。

至于正规性，还需额外验证（当 $F=\Q_p$ 时）：对可约 $(\varphi,\Gamma)$-模
\[
0\to\mathcal{R}_{F,L}(\delta_1)\to D\to\mathcal{R}_{F,L}(\delta_2)\to0
\]
的一般选取，叠 $\mathfrak{X}_2$ 在 $D$ 处光滑。这同样由切空间计算（相应地，Ext-群的计算）以及下述事实推出：在 $B$-丛的叠上一般地，比值 $\delta_1/\delta_2$ 非常正则，即
\[
\Ext^0_{\varphi,\Gamma}(\mathcal{R}_{F,L}(\delta_2),\mathcal{R}_{F,L}(\delta_1))=\Ext^2_{\varphi,\Gamma}(\mathcal{R}_{F,L}(\delta_2),\mathcal{R}_{F,L}(\delta_1))=0.
\]
\end{proof}

\begin{rem}{5.1.28}\label{rem:a4-5-1-28}
基本上同样的策略（不只使用 Borel 子群 $B$，而是使用一切抛物子群）应当能对一般的 $d$ 解决关于维数与局部结构的断言。只是扩张的空间/叠的维数计算变得更复杂，尚未完成。我们指出，Ext-群维数的类似计算也出现在 [BIP23a] 主要定理的证明中。
\end{rem}

我们简要阐述与第 4.1 节所述整理论之间的联系。设 $\Spf A\to\mathcal{X}_d$ 是一个 $A$-值点，其中 $\Z_p\langle T_1,\dots,T_m\rangle\twoheadrightarrow A$。

\ysrc{214}
该 $A$-值点定义了 $p$-进完备环 $\mathbf{A}_{F,A}=A\widehat{\otimes}_{\Z_p}\mathbf{A}_F$ 上的一个 $(\varphi,\Gamma)$-模 $D$（记号见注 5.1.4）；我们希望（函子性地）把 $\mathcal{R}_{F,A}[1/p]$ 上的一个 $(\varphi,\Gamma)$-模 $D_{\mathrm{rig}}$ 配给 $D$。由于不存在从 $\mathbf{A}_{F,A}$ 到 $\mathcal{R}_{F,A}[1/p]$ 的映射，必须先说明 $D$ 是超收敛的，即它容许子环 $\mathbf{A}^\dagger_{F,A}=A\widehat{\otimes}_{\Z_p}\mathbf{A}^\dagger_F$ 上的典范模型 $D^\dagger$（张量积取适当完备化）。若 $D$ 是附于 $A$ 上一族 $\mathrm{Gal}_F$-表示的 $(\varphi,\Gamma)$-模，这已由 Berger--Colmez 证明 [BC08, Cor. 4.2.7]。一般情形是 Gal Porat 的一个定理 [Por25]：

\begin{thm}{5.1.29}\label{thm:a4-5-1-29}
设 $A$ 是 $\Z_p$ 上拓扑有限型的 $p$-进完备 $\Z_p$-代数。则 $\mathbf{A}_{F,A}$ 上每个 étale $(\varphi,\Gamma)$-模都是超收敛的。
\end{thm}

作为该定理的推论，可以构造从第 4.1 节所定义叠 $\mathcal{X}_d$ 的刚性解析一般纤维到叠 $\mathfrak{X}_d$ 的态射。

\begin{cor}{5.1.30}\label{cor:a4-5-1-30}
存在态射
\begin{equation}
\pi_d:\mathcal{X}^{\mathrm{rig}}_{d,\eta}\longrightarrow\mathfrak{X}_d,\tag{5.1.31}
\end{equation}
它把 $(\varphi,\Gamma)$-模 $D\in\mathcal{X}_d(\Spf A)$（其中 $A$ 是某个拓扑有限型的 $p$-进完备 $\Z_p$-代数）送到 $D^\dagger\widehat{\otimes}_{\mathbf{A}^\dagger_{F,A}}\mathcal{R}_{F,A}[1/p]$。
\end{cor}

我们预期态射 $\pi_d$ 形式 étale，但不能由刚性解析空间（也不能由拟解析空间）表示。原因在于，该态射基本上是忘掉 Robba 环上 $(\varphi,\Gamma)$-模中的 étale 格，因而该态射的纤维参数化这类格的选取。在 $\mathrm{GL}_1$ 的情形，这意味着纤维形如叠商 $\GG_m/(\widehat{\GG}_m)^{\mathrm{rig}}$，其中 $\widehat{\GG}_m$ 是 $\Spf\Z_p$ 上的乘法形式群，细节见 §7.1。

\begin{rem}{5.1.32}\label{rem:a4-5-1-32}
我们指出，定理 5.1.20 的证明只给出刚性解析点局部附近的卡。由于用容许开子集对刚性解析空间所作的集合论覆盖未必是容许覆盖，这不足以给出猜想 5.1.19 的完整证明。把上述证明中的论证推广到相应 adic 空间的一切点（而不只是那些使 $L'=k(x)$ 在 $L$ 上有限的点 $x$），就能证明猜想 5.1.19。只要对 $\mathcal{R}_{F,L'}$ 上的 $(\varphi,\Gamma)$-模有 Kedlaya 斜率滤链定理的类比物（其中 $L'$ 可取为 $\Q_p$ 上有限型 adic 空间中任一点的（完备）剩余域），并且还能证明态射 $\pi_d$（它由 [Por25] 而存在）形式 étale，上述论证便可直接推广。étale 性预计不难证明，但 Kedlaya 定理向一般系数域的推广是否成立，我们毫无头绪。
\end{rem}

\subsubsection*{5.1.33. Hodge--Tate--Sen 权与整体截面}
$\mathrm{Gal}_F$ 在有限维 $\Q_p$-向量空间上的表示的 Hodge--Tate--Sen 权理论，可推广到 $(\varphi,\Gamma)$-模的族，相应地推广到 $\overline{X}_F$ 上 $\mathrm{Gal}_F$-等变向量丛的族。设 $A$ 是仿射代数，$D_A$ 是 $\mathcal{R}_{F,A}$ 上的 $(\varphi,\Gamma)$-模。记 $\mathcal{E}_A=\mathcal{E}_{F_\infty}(D_A)$ 为 $X_{F_\infty,A}$ 上相应的 $\Gamma$-等变向量丛。我们于是考虑特殊纤维
\[
\widehat{D}_{\mathrm{Sen}}(D_A)=\mathcal{E}_A|_{\Sp A\times\{\infty\}}=D_A\widehat{\otimes}_{\mathcal{R}_{F,A}}\bigl(\widehat{F}_\infty\widehat{\otimes}_{\Q_p}A\bigr).
\]

% 本片止于印刷页 214 页末（上式之后语句未完），续片从 PDF 228 页起
